\section{第三条线：从线性启动到爆发、次级磁岛与关停}\label{story:onset}
本线把时间顺序放回讨论中心。惯性两场模型先发现电流子层的持续变薄和非线性加速；暖电子、热离子及更高分辨率计算再揭示细丝、双重加速、次级岛与间歇减速。它们回答的是一个已被扰动的电流片如何发展。双稳态与自发转快研究则追问：系统为什么在某个时刻切换分支，且为什么同一参数会记得不同历史？

接着把线性Hall内层、磁岛切短电流片、递归破裂和非线性状态选择连起来。这个顺序十分重要：Hall改变初级增长、Hall在次级片中变强、开放层是否继续生岛，是不同阶段的三个判断。最后回到压强关停与能量转换，检验“更多岛、更快增长”是否真的对应更大的净磁通处理和能量释放。

\subsection{电子惯性重联：电流子层比电子惯性长度更薄}
\citet{Ottaviani1993}研究无耗散、不可压的二维冷两场模型，使用大 $\Delta'd_e$ 撕裂序列。其新发现不是简单的“电子惯性能断线”，而是\textbf{不同物理量出现不同层宽}：早期非线性阶段，速度与磁通扰动仍跨约 $d_e$ 的尺度，电流却产生不断缩窄的子层。广义磁通 $F$ 的拓扑仍被输运保存，普通磁通已改变。

论文用近似保持形状的宏观对流胞构造解析解。若 $A(t)$ 为流体位移幅度，在 $d_e\ll A\ll L$ 区间，子层厚度近似为 $d_e\exp[-A/\widehat d_e]$，其中 $\widehat d_e$ 与 $d_e$ 同阶；无量纲幅度方程含 $A''\sim A+cA^4$，体现非线性加速。这里的“准爆发”是有限适用区间内的加速，不能解释成真实等离子体产生无限电流。

本地JET预印本采用 $1024\times64$ Fourier分量，$L_x/L_y=0.5$、$d_e/(2L_x)=0.04$，仅一个线性不稳定模。作者在归一化时间约125时因子层无法解析而停止计算，并指出需要FLR、压力梯度、三维及可能的微不稳定性。这篇文章没有完成最终饱和和热化问题。

\subsection{暖电子约化模型：宏观饱和与微细结构继续演化共存}
\citet{Schep1994}构造的二维约化系统保存两族广义通量组合；其X点与普通磁通X点可以分离。有限电子质量使电流本身不再在相应定常传播解中奇异，但不保证电流梯度有界。该文的全FLR离子响应明确是满足若干极限与守恒要求的构造，而不是从Vlasov严格推出的任意尺度闭合；完整非线性本征解也没有求出。

\citet{L049}展示了这一结构的非线性后果。对于有限 $\rho_{\rm sound}/d_e$，$G_+$ 与 $G_-$ 在磁岛中沿不同方向卷绕，逐渐变成细丝。总能量仍守恒，磁能被转到小尺度流动和电子压缩等通道，粗粒化后的磁岛却可近似饱和。\textbf{“岛宽不再长大”与“所有自由度都停止演化”因此并不等价。}

冷电子极限则出现不同情况：初期电流更像条状层，随后速度剪切激发二次Kelvin--Helmholtz型不稳定性，产生单极与偶极涡、岛形脉动以及流体动能与电子动能间的振荡交换。随着 $\rho_{\rm sound}/d_e$ 增加到约1，所展示算例在岛饱和前不再出现该二次不稳定性。原文图1--3可直接比较暖电子细丝化、冷电子破碎和涡量结构。

这些图像不是任意参数下的相图。作者明确指出：无耗散暖电子层不断变薄，最终会使流体近似失效。随后用漂移动理学的正则动量薄片说明更丰富的相混合，但二维沿电流方向的对称性仍是关键条件。

\subsection{热离子改变细结构，未必明显改变总重联率}
Grasso等的三场陀螺流体模拟给出一个容易被只看磁通曲线漏掉的结果：热离子与热电子算例可有相近岛宽和发展速度，但静电势是否形成细丝、岛内密度是否变平、平行电场是否出现小尺度结构明显不同。热离子使静电势包含未经平滑的密度贡献，而陀螺平均势仍平滑；第一章已逐步推导这一关系。论文据此讨论流体电子束及可能增强的加速，但没有粒子分布函数，不能称为非热粒子谱的预测。\citep{E469}

另一个三维陀螺流体研究在小$\Delta'$参数下得到非线性爆发加速。最高扫描$\beta_e=1.6\times10^{-3}$时，以$L_\perp/c_s$归一的非线性峰值增长率为0.009，线性值0.0008。这里“增长率”来自扰动能量的对数导数，不能与$E/(B\VA)$直接比较。扫描同时改变盒子和初始片宽；X点分裂为两个Y点等拓扑变化在这篇短文中主要由文字报告，两幅结果图仅显示能量与增长率。三维是否必要、高速流是否是加速原因，作者仍留作后续问题。\citep{L184}

\subsection{爆发律的更新伴随计算分辨率的更新}
\citet{L191}以能量变分纠正冷惯性模型的非线性爆发指数，\citet{L214}在暖双场模型中强调局域X几何及$\dot\epsilon\propto\epsilon^{7/4}$。
\citet{F697}进一步比较离子/电子温度与弱电阻：固定岛幅时某些速率和能量释放曲线对微观尺度较弱敏感，但从种子开始的总时间仍依赖线性内层尺度。
更关键的是，其冷模型会出次级岛并间歇减速；作者明确指出早期横向模数少和数值黏性漏掉了磁岛。
因此新结果是“局域爆发与磁岛间歇相竞争”，而不是一条适用于所有无碰撞重联的幂律。

\subsection{陀螺诱导双重加速：先出现细尺度强流，再打开宏观排气}
\citet{F672}用强导引场四场陀螺流体，在大$\Delta'$、热离子和微观尺度远小于系统尺度时看到：线性增长之后先急剧加速，继而短暂减速或停滞，随后再次加速进入快重联。原文图6、8、9、11应连着看，分别追踪增长史、静电势、流场和排气开口；只看峰值会把两个过程合并。

第一轮加速与热离子有限回旋半径有关。在保持组合热尺度$\rho_\tau$相同、互换热离子／热电子贡献的对照中，热离子例出现更强的局地电势细结构，$E\times B$流可约高一个数量级；冷离子例缺少同样的第一轮加速。第一章已从准中性算子说明：热离子内区电势可直接跟随细密度结构，冷离子内区则经逆Laplacian平滑。这是与可测场量相连的机制对照。

第二轮加速随排气结构进一步开放，冷热离子均可出现。小尺度与宏观尺度分离越明显，两轮之间的减速越清楚。原文所画$\gamma_{\rm eff}=\dd\ln\delta\psi/\dd t$是\emph{对数磁通增长率}；它的尖峰不等于重联电场本身的尖峰，不能拿来与0.1直接比较。

表I将$d_e$从0.2减至0.025时，归一化峰值率由约0.283逐渐接近0.100；分母采用\emph{初始最大面内磁场}，不是随时间测得的局地上游场。作者也限定$d_e\Delta'\gtrsim1$：若固定$\Delta'$继续令$d_e\to0$，最终会进入小岛饱和区。因此该表不能证明严格理想极限仍有普适有限率。谱滤波、最细结构失去解析后的停止时刻和电子回旋效应缺失，同样限制了后期判断。

\subsubsection{哪些部分已经有动理学对照，哪些还没有}
\citet{E475}比较四场陀螺流体与EUTERPE陀螺动理学，先让两种代码各自对照其模型的线性本征解，再进行跨模型比较。这能把“离散是否解对自身方程”与“两个模型是否近似同一物理”分开。

在所扫线性波数中，增长率差别可达约20\%；减小平衡面内场后差别减弱，但没有自动降到零。非线性部分主要检验\emph{小}$\Delta'$：两模型都得到小岛饱和及其振荡，振荡频率很接近；饱和半宽差在较大驱动的$k_y=1.6$处约30\%，接近边缘稳定时明显减小。岛宽对温度相对不敏感，振荡频率却可随离子温度增加而近乎加倍。

论文文字称一个温比扫描的岛宽仅变约5\%，但图11显示的变化明显更大；据图估读，动理学约从0.040到0.049，流体约从0.046到0.053。这里只把它判为原文内部数量不一致，不把图上估读当作精确源数据。更关键的是，动理学非线性计算保留的模主要为$m=0,\pm1$，参数也在小岛区；\textbf{它没有检验前文大$\Delta'$的两次加速}。振荡可与粒子反弹频率比较，不代表有限矩模型恢复了完整俘获分布。

\subsection{Full-F、任意波长极化与多X点不同步}
GREENY及其后续论文\citep{E473,L314}显示，保留任意波长极化后，某些算例比长波近似发展更快并出现更细的电势和电流结构；全密度模型还保留了小振幅模型忽略的乘积项。初始电流由电子还是离子承担会明显改变过程，后续研究中离子载流启动通常慢约一个数量级。拉长盒子后可有多个X点在不同阶段同时活动，不能用一条单X点曲线代表整个域。

非正规线性算子提供暂态放大的可能，非线性阶段又伴随谱展宽；两者不能直接画等号。暖离子会改变态的形貌，但主岛数不随温比稳定单调增加。若只看摘要中的“新模型更快”或“非正规性导致爆发”，容易漏掉正文对范数、排序、初值和参数标签的限制。第一章给出了可以自行检查的能量范数和极化推导。

\subsection{快慢双稳态、迟滞与突发转换}
\citet{Cassak2005}的重要证据是同一参数、不同启动历史下出现两种状态。在其F3D算例中，归一化电阻约0.015时可得约0.06的快支与约0.01的慢支；降低电阻后再恢复，系统可留在快支。原文图3、4给出两个方向不同的转换阈值，这比仅比较两个不同参数算例更有力。

同时应记录，电子惯性长度为0.2而网格为0.1，只有两格跨一个 $d_e$；2005年正文没有单独展示网格收敛。对DIII-D的温度估算低约一个数量级，MRX比较当时缺少片宽数据。因此模拟中的迟滞证据与实验的定量验证必须分别评价。

\citet{Cassak2010}给出鞍结分岔的简化动力学解释，并在关闭电子惯性的Hall计算中仍得到双稳态。但它保留电阻和四阶耗散，并非无耗散证明；描述片宽变化的方程还没有自洽决定局地长度和演化中的上游场。该文对电阻平衡Hall态稳定性的论断，与后文的大系统结果存在需要区分参数和几何的争议。

\subsection{强场被带入后的自发启动：保持电阻也能转快}
Cassak等2006年从宽正弦片出发，发现慢重联把越来越强的上游场输送到内嵌SP层。电阻保持不变时，层仍逐渐变薄；达到当地$d_i$附近后离子电子解耦，出现开放排气与快速流入。慢阶段符合$\eta/\delta$，快阶段采用$0.17V_{\!A,\rm up}$拟合。这给出“先积累、再迅速释放”的一种具体模型过程。\citep{Cassak2006}

展示的低电阻运行早期由另一电阻运行重启，启动前瞬变已消失；因此准确说法是后期转变不需要在触发时继续改变电阻。太阳参数下的积累及转换时间只是量级相容，既未直接解析太阳的电子层，也未纳入高$S$磁岛、热反馈及三维竞争，不能据此唯一判断耀斑启动机制。

\subsection{线性Hall加速有条件，不能当成理想Hall直接断线}
\citet{L123}研究线绑边界的线性X点衰减；弱Hall拟合同时依赖$d_i$和电阻，因此将尺寸变化只归到$d_i^2$会漏掉归一化电阻的变化。
\citet{L122}提出三层变分近似，却没有完整$\Delta'$外匹配，且高斯试探电流导致某个$1/k^2$积分需额外正则化说明。其弱Hall修正还存在可直接代数复核的系数问题。
故本报告保留其启发，不将其形式上有限的无电阻增长解释成平滑理想Hall中磁通已被打破。

更系统的Hall撕裂分析\citep{L250}区分可压缩、常磁通与非常磁通分支；某一分支增长不显含$S$，仍需有限电阻完成内层匹配。
\citet{L238}在受迫稳定Taylor平衡中找到尤其值得注意的中间区：早期已出现Hall动力学，后期却能回到MHD式受限演化，尽管片厚小于$d_i$。
Hall还可能提高单模增长却增大最快波长，使可容纳的磁岛更少；“更快增长”与“更多磁岛”不是同一句话。

\subsection{磁岛先改变内层，再触发Hall阶段：中间态可以往复切换}
\citet{Baalrud2011}分清了两个阈值。长薄片的整体宽度可以仍大于 $d_i$，但撕裂扰动的电阻内层 $d_\eta$ 已小于 $d_i$，Hall项便可改变增长率与本征结构。适用排序为 $d_\eta\ll d_i\ll\delta_{\rm SP}$，不是整体电流片已经成为Hall开口层。

在该排序、快速增长和长波近似下，理论得到
\begin{equation}
{\gamma_{\max}\over\Gamma_0}\simeq3.2(d_i/L)^{6/13}S^{7/13},\quad
k_{\max}L\simeq2.9(d_i/L)^{1/13}S^{11/26},\quad\Gamma_0=2\VA/L.
\end{equation}
全非线性计算提取的初期增长率，在 $S=10^6$、$d_i/L\simeq2\times10^{-4}$ 附近从弱依赖变为明显依赖Hall尺度，与预测转折接近；作者也明确说，有限参数区和增长率提取误差不足以精确验证 $6/13$ 指数。图8显示长片中的小磁岛已带四极面外场。

\citet{Huang2011}研究另一层次：磁岛把长片切成短片，使局地片先达到Hall尺度。$S=5\times10^5$ 的扫描中，$L/d_i=2500$ 可保持开放单X点，5000时却在开放层、延长层和多磁岛状态之间循环，速率约在0.005--0.03变化。10000的情况没有进入相同Hall阶段。\textbf{“进入Hall尺度”不保证最终永久停在单X点。}

这些比较还受厚度定义影响：原文某时刻四分之一峰值全宽约 $0.7d_i$，半高半宽却约 $0.17d_i$。用不同定义讨论“片宽是否超过 $d_i$”会造成假矛盾。该文还说明，自然系统的尺度参数可能远离其循环中间区，不应把循环当作所有耀斑或磁尾的必然模式。

\subsection{递归片破裂：初级Hall很弱，次级Hall仍能显著增强}
Pucci等线性本征计算显示，Hall作用可以在$d_i<a$却$d_i>\delta_\eta$时改变撕裂，形成变宽的Hall中间层和变薄的内电阻层。拟合的增长指数与简单量纲匹配不同，并给出修正的临界长宽比；推导和适用范围已集中在第一章。\citep{L236}

Papini等的非线性计算进一步发现，初级磁岛之间自发形成半厚约初始$0.1$--$0.15$的次级片，片可先保持近压强平衡，再迅速破裂。按原片内层计$d_i/\delta_\eta=0.2$时，Hall几乎不改变初级增长，却使次级增长增加约17\%；比值1.4时次级增长接近加倍。关键新情况是：一个在宏观初级阶段可忽略的项，经由递归尺度缩小成为重要项。\citep{L262}

同时必须保留三条未决限制。第一，文中$\gamma_{\rm rec}$是对数磁通增长或磁扰动指数拟合，并非通常的$E/(B\VA)$；次级$2.3$与初级$0.48$不能与“普适0.1”直接比较。第二，换成各次级片自己的阿尔芬时间后，增长仍为一的量级；实测Hall次级片长宽比却\emph{不吻合}前述线性拟合标度，原文明确承认参数扫描不足。第三，强Hall下次级片不再可见，可能是更快破裂或X点几何改变，也可能是分辨率不足，不能写成已确证的“Hall必然抑制磁岛”。

\subsection{递归破裂后的三种状态：开放X点仍可持续排出磁岛}
Shi等在固定电阻、相同随机种子的七组Harris片计算中，观察到随$d_i$增强发生的变化：纯MHD反复形成更短电流片；弱Hall在开放X点继续产生并排出次级岛；较强Hall则形成不再明显生岛的开放层。排气开口角约由10度增至28度。这个中间态不是“Hall态与长SP片反复跳转”，而是开放结构与持续生岛共存。\citep{L263}

磁通斜率从约0.032升至0.058，但MHD晚期递归破裂后的0.034已经超过最弱Hall算例。较强Hall层长逐渐不变时，其半长仅剩约3--4个横向网格；最小MHD片也达到网格尺度。因此，观测到的状态差异比“已证实普适0.06上限”更可靠。它与Papini等的次级增长增强并不矛盾：线性增长、非线性开口和后期是否继续生岛，是三个阶段不同的问题。\citep{L262,L263}

\subsection{次级磁岛并不总带来更快的净能量释放}
三个证据应并列而非互相替代。第一，磁岛合并的强导引场PIC会不断生成次级岛，Hall却维持开放X形，二者平均率仍相近\citep{L241}。第二，PIC质量比扫描中低质量比容易生成岛，但岛并未使全局能量转换像电阻MHD算例那样显著提高；主电子层和移动次级X点的非理想预算也不同\citep{F606}。第三，欠分辨的MHD可以产生额外岛和假平台率\citep{F788}。

因此值得研究的新问题是：在什么驱动、寿命和尺度区间，岛的生成真正控制磁通通量？不能仅用“有岛”推断“岛介导的率”，也不能用一个无岛的已分辨算例否定其他初值下的撕裂。质量比趋于真实值、离子FLR、导引场、噪声和开放边界仍需联动检验。

\subsection{EMHD岛内部还会发生另一种不稳定性}
\citet{E459}在冷惯性EMHD中发现，早期有序的电流片级联被岛内射流的Kelvin--Helmholtz类涡旋破坏；电子压缩性可延迟或抑制这一过程。
它考察的主要是磁岛内部电流/涡量重新组织，不应只按“撕裂产生更多磁岛”概括。
模拟最高另做到$2048\times512$；晚期能量/Casimir误差可达5\%--10\%，不能由衰减曲线直接推断真实电子加热。

\citet{L202}则在惯性/电阻EMHD线性问题中强调纯增长本征模与驻波whistler不同，并构造有限中心电流的解。
其与MRX的增长时间和宽度只是尺度估计：计算约1.2 cm宽、实验比较约3.3 cm，且从假定10\%初扰动估总时间；未直接模拟稳态MRX排出和全欧姆预算。

\citet{L230}的强导场X点塌缩结果来自PIC，应放在动理学对照类。
其中电子涡旋可伴随次级X点，部分区域由对流惯性支撑电场；三维出现倾斜涡管。
但其三维图没有相应完整欧姆预算或全局率诊断，$Q$涡旋指标也不是重联判据，不能据图证明“电子压力项在三维全处消失”。

\subsection{磁岛触发Hall之后，重联还可能因压强结构被关停}
\citet{L159}在一组大$S$计算中看到Sweet--Parker、磁岛增强、Hall快速阶段先后出现，Hall阶段排出竞争磁岛。
后续文献已有开放X点持续排岛的情况，因此“Hall与磁岛互斥”只能限于该组边界、尺度与时段。
\citet{L174}则把压强梯度带入重联层，发现即使还有可释放磁能，反向抗磁漂移相对出流足够大时也可抑制重联。
该模型强导场、冷离子、绝热电子压力，MAST比较又依赖尺度估计；它提出了饱和/关停候选机制，尚非完整托卡马克锯齿预测。

\citet{E472}研究的是另一个问题：初始密度梯度使线性岛传播并降低增长；热离子时传播可朝离子抗磁方向。
其解析频率/增长率与数值有可见偏差，并未等价地验证\citet{L174}的非线性出流竞争条件。

\paragraph{这一线索走到了哪里。}已发现的状态远多于一个稳定X点：冷层子尺度尖化、暖模型细丝化、两次加速、快慢迟滞、反复长片与开放层、开放X点持续排岛、次级剪切涡，以及磁能尚存时的压强关停。尚缺的是在同一初边值、充分分辨和统一率定义下连接这些状态的相图。任何新出现的磁岛或爆发指数，都应先通过层宽、谱尾、噪声和寿命检查；第2章的数值检验因此是这条发现链的一部分。
